\subsection{Sets and Proper Classes}

\begin{proposition}[Proper Classes Are Never Members]\label{prop:proper-classes-are-never-members}
If \(A\) is a proper class, then \(A\) is not a member of any class \(B\).

In symbols:
\[\forall A \forall B (\PC(A) \to A \notin B)\]
\end{proposition} \vspace{0.3cm}

\begin{proof}
Let \(A\) and \(B\) be arbitrary classes, and suppose that \(A\) is a proper class. Assume, for a contradiction, that \(A\in B\). Then \(A\) must be a set, contradicting the definition of a proper class. Hence \(A\notin B\). Discharging the assumption that \(A\) is a proper class yields \(\PC(A)\to A\notin B\). Since \(A\) and \(B\) were arbitrary, universal generalization establishes the proposition.
\end{proof}

\begin{formalderivation}
\begin{prooftree}
    \AxiomC{}
    \RightLabel{\(1\)}
    \UnaryInfC{\(A\in B\)}
    \RightLabel{Set-I}
    \UnaryInfC{\(\Set(A)\)}
%
    \AxiomC{}
    \RightLabel{\(2\)}
    \UnaryInfC{\(\PC(A)\)}
    \RightLabel{\(\mathrm{Def.}\)}
    \dashedLine
    \UnaryInfC{\(\neg\Set(A)\)}
%
    \RightLabel{\(\neg\)-E}
    \BinaryInfC{\(\bot\)}
    \RightLabel{\(\neg\)-I, \(1\)}
    \UnaryInfC{\(A \notin B\)}
    \RightLabel{\(\to\)-I, \(2\)}
    \UnaryInfC{\(\PC(A)\to A \notin B\)}
    \RightLabel{$\forall$-I*}
    \UnaryInfC{$\forall A \forall B(\PC(A) \to A \notin B)$}
\end{prooftree}
\end{formalderivation} \bookvspace{0.5cm}

\begin{proposition}[Subclasshood Is Reflexive]\label{prop:subclasshood-reflexive}
\[
\forall x(x\subseteq x).
\]
\end{proposition}

\begin{proof}
Let \(A\) be an arbitrary class. To show that \(A \subseteq A\), we must show that every element of \(A\) is also an element of \(A\). But this is immediate: if \(x \in A\), then \(x \in A\). Hence \(A \subseteq A\). Since $A$ was arbitrary, Proposition \ref{prop:subclasshood-reflexive} follows by Universal Generalization.
\end{proof}

\begin{formalderivation}
    \begin{prooftree}
    \AxiomC{}
    \RightLabel{\(1\)}
    \UnaryInfC{\(a\in A\)}
    \RightLabel{\(\subseteq\)-I, \(1\)}
    \UnaryInfC{\(A\subseteq A\)}
\end{prooftree}
\end{formalderivation}